% !TEX spellcheck = en_US



% ~~~~~~~~~~~~~~~~~~~~~~~~~~~~~~~~~~~~~~~ V
% (NC) 2026 Haluk Bingol
% github.com/halukbingol/LaTeX-Templates
%
% Licensees may copy, distribute, display, and perform the work and make derivative works 
% and remixes based on it only for non-commercial purposes. 
% ~~~~~~~~~~~~~~~~~~~~~~~~~~~~~~~~~~~~~~~ A




% v20261030T1200

\documentclass[aspectratio=169,10pt,compress]{beamer}




% ~~~~~~~~~~~~~~~~~~~~~~~~~~~~~~~~~~~~~~~~~~~~~~~~~~~~~~~~~~~~~~~~~ V hblib
	\usepackage{../../hblib/hblibPresentation} % lib presentation: theme, i/N, \hSection, algorithm2e, ...
	\usepackage{../../hblib/hblibCodeoutput} % lib code + output: \lstcode, \lstterminal, ...
	\usepackage{../../hblib/hblibDatastructures} % lib data structures: \hString, \hStack, \hQueue, \hArray, \hGrid
% ~~~~~~~~~~~~~~~~~~~~~~~~~~~~~~~~~~~~~~~~~~~~~~~~~~~~~~~~~~~~~~~~~ A hblib




% +++++++++++++++++++++++++++++++++++++++ V specific
\DontPrintSemicolon % in algorithms
% +++++++++++++++++++++++++++++++++++++++ A specific



\title[Arrays ND: Junior]{
	Multi-Dimensional Arrays\\
	Junior
}
\subtitle{
	\cisOOO\ -- Level 3: ragged structures, matrices, performance
}
\author[Erkovan]{Sait Erkovan}

\institute[FCIS]{
	Faculty of Computer and Information Sciences\\
	Yeditepe University
}
\date{{\tiny \hbTimeStamp}}

\begin{document}




% ~~~~~~~~~~~~~~~~~~~~~~~~~~~~~~~~~~~~~~ V frm
\begin{frame}[t,fragile]
	\titlepage
\end{frame}
% ~~~~~~~~~~~~~~~~~~~~~~~~~~~~~~~~~~~~~~ A frm




% ~~~~~~~~~~~~~~~~~~~~~~~~~~~~~~~~~~~~~~ V frm
\begin{frame}[t,fragile] \frametitle{Outline}
	\begin{columns}[T]
		\begin{column}{0.48\textwidth}
			\tableofcontents[sections={1-3}, hideallsubsections]
		\end{column}
		\begin{column}{0.48\textwidth}
			\tableofcontents[sections={4-5}, hideallsubsections]
		\end{column}
	\end{columns}

	\vspace{1em}
	\small It builds on \hCode{array1D-3-junior}: costs, prefix sums and locality in 1D are assumed.
\end{frame}
% ~~~~~~~~~~~~~~~~~~~~~~~~~~~~~~~~~~~~~~ A frm




% ~~~~~~~~~~~~~~~~~~~~~~~~~~~~~~~~~~~~~~ V
\hSection[Ragged]{Ragged structures}
% ~~~~~~~~~~~~~~~~~~~~~~~~~~~~~~~~~~~~~~ A




% ~~~~~~~~~~~~~~~~~~~~~~~~~~~~~~~~~~~~~~ V frm
\begin{frame}[t,fragile] \frametitle{Pascal's triangle}
	\lstcoderightstep[0.66][basicstyle=\ttfamily\scriptsize, firstline=3, lastline=18, firstnumber=3]{codeoutput/Pascal.java}{Java}{codeoutput/Pascal.txt}

	\footnotesize Row \hCode{r} has \hCode{r + 1} cells, so a ragged table fits exactly: no unused cells.
	Every inner cell is the sum of the two cells above it.
\end{frame}
% ~~~~~~~~~~~~~~~~~~~~~~~~~~~~~~~~~~~~~~ A frm




% ~~~~~~~~~~~~~~~~~~~~~~~~~~~~~~~~~~~~~~ V frm
\begin{frame}[t,fragile] \frametitle{A triangular table saves half}
	\lstcoderightstep[0.66][basicstyle=\ttfamily\scriptsize]{codeoutput/TriangleTable.java}{Java}{codeoutput/TriangleTable.txt}

	\footnotesize A distance table is symmetric: $d(a, b) = d(b, a)$. Storing only $c \le r$ needs
	$n(n + 1)/2$ cells instead of $n^{2}$; \hCode{get} hides the order of the two indices.
\end{frame}
% ~~~~~~~~~~~~~~~~~~~~~~~~~~~~~~~~~~~~~~ A frm




% ~~~~~~~~~~~~~~~~~~~~~~~~~~~~~~~~~~~~~~ V
\hSection[Matrices]{Matrices}
% ~~~~~~~~~~~~~~~~~~~~~~~~~~~~~~~~~~~~~~ A




% ~~~~~~~~~~~~~~~~~~~~~~~~~~~~~~~~~~~~~~ V frm
\begin{frame}[t,fragile] \frametitle{The matrix product: a row times a column}
	\begin{center}
		\begin{tikzpicture}
			\hGrid[id=a, name={a}, width=0.55, rowhot={0}]{{1,2,3},{4,5,6}}
			\node at (2.15,-0.2) {$\times$};
			\hGrid[id=b, at={(2.9,0.25)}, name={b}, width=0.55, colhot={0}]{{1,0},{0,1},{2,2}}
			\node at (4.4,-0.2) {$=$};
			\hGrid[id=c, at={(5.0,0)}, name={c}, width=0.55, hot={0/0}]{{7,8},{16,17}}
		\end{tikzpicture}
	\end{center}
	\[
		c[i][j] = \sum_{x=0}^{k-1} a[i][x] \cdot b[x][j], \qquad c[0][0] = 1 \cdot 1 + 2 \cdot 0 + 3 \cdot 2 = 7
	\]
	\begin{itemize}\small
		\item $a$ is $n \times k$, $b$ is $k \times m$: the inner sizes must agree; $c$ is $n \times m$
		\item Every cell of $c$ is a sum of $k$ products: $n \cdot m \cdot k$ multiplications, three nested loops
		\item The order matters: $a \cdot b \ne b \cdot a$ in general
	\end{itemize}
\end{frame}
% ~~~~~~~~~~~~~~~~~~~~~~~~~~~~~~~~~~~~~~ A frm




% ~~~~~~~~~~~~~~~~~~~~~~~~~~~~~~~~~~~~~~ V frm
\begin{frame}[t,fragile] \frametitle{The product in Java}
	\lstcoderightstep[0.66][basicstyle=\ttfamily\scriptsize, firstline=3, lastline=20, firstnumber=3]{codeoutput/MatMul.java}{Java}{codeoutput/MatMul.txt}

	\footnotesize The sizes come from the arrays: \hCode{a.length} rows, \hCode{b.length} $= k$, \hCode{b[0].length} columns.
\end{frame}
% ~~~~~~~~~~~~~~~~~~~~~~~~~~~~~~~~~~~~~~ A frm




% ~~~~~~~~~~~~~~~~~~~~~~~~~~~~~~~~~~~~~~ V frm
\begin{frame}[t,fragile] \frametitle{Three loops: $n^{3}$, measured}
	\lstcoderightstep[0.66][basicstyle=\ttfamily\scriptsize, firstline=23, lastline=32, firstnumber=23]{codeoutput/MatMulGrowth.java}{Java}{codeoutput/MatMulGrowth.txt}

	\begin{itemize}\footnotesize
		\item Doubling $n$ multiplies the work by 8; the times grow by about that (or more, when the tables leave the cache)
		\item Square matrices of size 10\,000 would need $10^{12}$ steps: libraries use better orders and many cores
		\item The times differ from run to run and from machine to machine
	\end{itemize}
\end{frame}
% ~~~~~~~~~~~~~~~~~~~~~~~~~~~~~~~~~~~~~~ A frm




% ~~~~~~~~~~~~~~~~~~~~~~~~~~~~~~~~~~~~~~ V frm
\begin{frame}[t,fragile] \frametitle{Transpose a square table in place}
	\lstcoderightstep[0.62][basicstyle=\ttfamily\scriptsize]{codeoutput/TransposeInPlace.java}{Java}{codeoutput/TransposeInPlace.txt}

	\footnotesize Swap $(r, c)$ with $(c, r)$ only above the diagonal: $c$ starts at $r + 1$.
	Starting at 0 would swap every pair twice and give back the original.
\end{frame}
% ~~~~~~~~~~~~~~~~~~~~~~~~~~~~~~~~~~~~~~ A frm




% ~~~~~~~~~~~~~~~~~~~~~~~~~~~~~~~~~~~~~~ V frm
\begin{frame}[t,fragile] \frametitle{Rotate a picture by 90 degrees}
	\lstcoderightstep[0.66][basicstyle=\ttfamily\scriptsize]{codeoutput/Rotate.java}{Java}{codeoutput/Rotate.txt}

	\footnotesize Row \hCode{r} of the picture becomes column \hCode{rows - 1 - r} of the result:
	a transpose followed by a mirror. A $2 \times 3$ picture becomes $3 \times 2$, so a new table is needed.
\end{frame}
% ~~~~~~~~~~~~~~~~~~~~~~~~~~~~~~~~~~~~~~ A frm




% ~~~~~~~~~~~~~~~~~~~~~~~~~~~~~~~~~~~~~~ V
\hSection[Performance]{Performance}
% ~~~~~~~~~~~~~~~~~~~~~~~~~~~~~~~~~~~~~~ A




% ~~~~~~~~~~~~~~~~~~~~~~~~~~~~~~~~~~~~~~ V frm
\begin{frame}[t,fragile] \frametitle{Row by row or column by column? Measured}
	\lstcodedownstep[basicstyle=\ttfamily\fontsize{6.4}{7.6}\selectfont, firstline=5, lastline=19, firstnumber=5][basicstyle=\ttfamily\scriptsize]{codeoutput/LoopOrder.java}{Java}{codeoutput/LoopOrder.txt}

	\footnotesize \uncover<2->{The same $16 \cdot 10^{6}$ additions; the column order is many times slower.}
\end{frame}
% ~~~~~~~~~~~~~~~~~~~~~~~~~~~~~~~~~~~~~~ A frm




% ~~~~~~~~~~~~~~~~~~~~~~~~~~~~~~~~~~~~~~ V frm
\begin{frame}[t,fragile] \frametitle{Why: a column jumps from row to row}
	\begin{center}
		\begin{tikzpicture}
			\hArray[id=o, at={(0,-0.35)}, width=0.5, index=none, name={t}]{\phantom{0},\phantom{0},\phantom{0}}
			\hArray[id=r0, at={(2.6,0.9)}, width=0.5, index=none, done={0,...,5}]{,,,,,}
			\hArray[id=r1, at={(4.0,-0.1)}, width=0.5, index=none, cur={0}]{,,,,,}
			\hArray[id=r2, at={(2.0,-1.1)}, width=0.5, index=none, cur={0}]{,,,,,}
			\draw[href] (o-0.center) -- (r0-0.west);
			\draw[href] (o-1.center) -- (r1-0.west);
			\draw[href] (o-2.center) -- (r2-0.west);
			\node[font=\scriptsize, anchor=west] at (5.8,1.15) {row 0: read in one sweep};
			\node[font=\scriptsize, anchor=west] at (7.2,0.15) {column 0: one cell here,};
			\node[font=\scriptsize, anchor=west] at (5.2,-0.85) {one cell there, \dots};
		\end{tikzpicture}
	\end{center}

	\begin{itemize}\small
		\item The cells of one row are side by side in memory; the rows are separate arrays, somewhere else
		\item Row by row reads memory in order: the cache and the processor love it (\hCode{array1D-3-junior})
		\item Column by column touches a different row array at every step: each access may miss the cache
		\item Rule: make the \emph{last} index the inner loop
	\end{itemize}
\end{frame}
% ~~~~~~~~~~~~~~~~~~~~~~~~~~~~~~~~~~~~~~ A frm




% ~~~~~~~~~~~~~~~~~~~~~~~~~~~~~~~~~~~~~~ V frm
\begin{frame}[t,fragile] \frametitle{One flat array instead of a table}
	\lstcoderightstep[0.66][basicstyle=\ttfamily\scriptsize]{codeoutput/FlatTable.java}{Java}{codeoutput/FlatTable.txt}

	\begin{itemize}\footnotesize
		\item Cell $(r, c)$ is \hCode{flat[r * cols + c]}: row-major order by hand
		\item One object, no row references, all cells side by side: images, games and numeric libraries do this
		\item The price: no ragged rows, no \hCode{t[r]}, and the index formula must be right everywhere
	\end{itemize}
\end{frame}
% ~~~~~~~~~~~~~~~~~~~~~~~~~~~~~~~~~~~~~~ A frm




% ~~~~~~~~~~~~~~~~~~~~~~~~~~~~~~~~~~~~~~ V
\hSection[Techniques]{Techniques}
% ~~~~~~~~~~~~~~~~~~~~~~~~~~~~~~~~~~~~~~ A




% ~~~~~~~~~~~~~~~~~~~~~~~~~~~~~~~~~~~~~~ V frm
\begin{frame}[t,fragile] \frametitle{Summed-area table: any rectangle in O(1)}
	\lstcodedownstep[basicstyle=\ttfamily\scriptsize, firstline=3, lastline=16, firstnumber=3][basicstyle=\ttfamily\scriptsize]{codeoutput/SummedArea.java}{Java}{codeoutput/SummedArea.txt}

	\footnotesize The prefix sums of \hCode{array1D-3-junior}, in two dimensions: one pass builds \hCode{s},
	then every rectangle sum is 4 lookups. One extra row and column of zeros remove the edge cases.
\end{frame}
% ~~~~~~~~~~~~~~~~~~~~~~~~~~~~~~~~~~~~~~ A frm




% ~~~~~~~~~~~~~~~~~~~~~~~~~~~~~~~~~~~~~~ V frm
\begin{frame}[t,fragile] \frametitle{Game of Life: neighbors and a second table}
	\lstcode[basicstyle=\ttfamily\fontsize{6.4}{7.6}\selectfont, firstline=2, lastline=19, firstnumber=2]{codeoutput/Life.java}{Java}

	\footnotesize A cell lives on with 2 or 3 live neighbors and is born with exactly 3. All cells change \emph{at once}:
	the new values go into \hCode{next}, never into \hCode{g} while it is still being read.
\end{frame}
% ~~~~~~~~~~~~~~~~~~~~~~~~~~~~~~~~~~~~~~ A frm




% ~~~~~~~~~~~~~~~~~~~~~~~~~~~~~~~~~~~~~~ V frm
\begin{frame}[t,fragile] \frametitle{Game of Life: three generations}
	\lstcoderightstep[0.66][basicstyle=\ttfamily\scriptsize, firstline=21, lastline=34, firstnumber=21]{codeoutput/Life.java}{Java}{codeoutput/Life.txt}

	\footnotesize \hCode{gens} is a three-dimensional array: generation, row, column. The blinker turns
	every step and comes back after two.
\end{frame}
% ~~~~~~~~~~~~~~~~~~~~~~~~~~~~~~~~~~~~~~ A frm




% ~~~~~~~~~~~~~~~~~~~~~~~~~~~~~~~~~~~~~~ V
\hSection[Summary]{Summary}
% ~~~~~~~~~~~~~~~~~~~~~~~~~~~~~~~~~~~~~~ A




% ~~~~~~~~~~~~~~~~~~~~~~~~~~~~~~~~~~~~~~ V frm
\begin{frame}[t,fragile] \frametitle{Cheat sheet: what each operation costs}


% +++++++++++++++++++++++++++++++++++++++ V tbl
%: -Tbl. costs
\begin{table}
		\small
		\begin{tabular}{@{}L{0.40\textwidth}L{0.16\textwidth}L{0.36\textwidth}@{}}
			\toprule
			\textbf{Operation, $n \times m$ table} & \textbf{Steps} & \textbf{Note} \\
			\midrule
			read or write one cell                  & $O(1)$          & two lookups: the row, the cell \\
			visit every cell                        & $O(n m)$        & last index inside: faster \\
			rectangle sum, summed-area table        & $O(1)$          & after $O(n m)$ preparation \\
			transpose, rotate                       & $O(n m)$        & a new table unless square \\
			product of $n \times k$ and $k \times m$ & $O(n k m)$      & $O(n^{3})$ for square tables \\
			one Life step                           & $O(n m)$        & 8 neighbors per cell, a second table \\
			triangular storage                      & $n(n+1)/2$ cells & instead of $n^{2}$ \\
			\bottomrule
		\end{tabular}
\end{table}%
% +++++++++++++++++++++++++++++++++++++++ A tbl
\end{frame}
% ~~~~~~~~~~~~~~~~~~~~~~~~~~~~~~~~~~~~~~ A frm




% ~~~~~~~~~~~~~~~~~~~~~~~~~~~~~~~~~~~~~~ V frm
\begin{frame}[t,fragile] \frametitle{Exercises and projects}
	\begin{enumerate}\small
		\item \hCode{MatMulGrowth}: swap the two inner loops ($i, x, j$). Measure again and explain the difference.
		\item Write \hCode{rotateLeft} and check that four left turns give back the original picture.
		\item \hCode{SummedArea}: count how many rectangles of a table have a sum above a limit, in $O(n^{2} m^{2})$.
		\item \hCode{TriangleTable}: store a symmetric $n \times n$ table in one flat array of $n(n+1)/2$ cells.
	\end{enumerate}

	\vspace{0.5em}
	\hRemark{Projects}: Game of Life with a ``wrap around'' board; a small image editor for \hCode{int[][]} gray
	pictures (blur, flip, rotate, histogram); a matrix class with \hCode{times}, \hCode{transpose}, \hCode{equals}.
\end{frame}
% ~~~~~~~~~~~~~~~~~~~~~~~~~~~~~~~~~~~~~~ A frm




% ~~~~~~~~~~~~~~~~~~~~~~~~~~~~~~~~~~~~~~ V frm
\begin{frame}[t,fragile] \frametitle{References}
	\small Matrix multiplication \cite{cormen2022introduction}; arrays in memory \cite{knuth1973artv1}.

	\vspace{1em}
	\setbeamertemplate{bibliography item}[text]
	\bibliographystyle{IEEEtran}
	\tiny\bibliography{../../hblib/hblibDataStructures}
\end{frame}
% ~~~~~~~~~~~~~~~~~~~~~~~~~~~~~~~~~~~~~~ A frm




% ~~~~~~~~~~~~~~~~~~~~~~~~~~~~~~~~~~~~~~ V frm
\begin{frame}[t,fragile]
	\begin{center}
		\vspace{4cm}
		\hRemark{\Huge Questions?}
	\end{center}
\end{frame}
% ~~~~~~~~~~~~~~~~~~~~~~~~~~~~~~~~~~~~~~ A frm



\end{document}
